% Transcription — fonds Alexandre Grothendieck, shelfmark 15, batch 10 (pages 181-200).
% Established from the facsimile archives/batches/15/batch-10.pdf (a mirror of
% https://grothendieck.umontpellier.fr/15.pdf), per the skill
% transcribe-grothendieck. The batch file carries no cover sheet: PDF page N
% is page 180+N of the folder (the Montpellier footer prints « 181 » ... « 200 »),
% and the archivists' pencil numbers bottom-left agree wherever checked (182,
% 184, 186, 190, 194, 196, 198, 200).
%
% Pass: Opus 5.5 (claude-opus-5-5), 2026-09-23 — first pass, unchecked against the pages by a human.
%
% Ten of the twenty pages are transcribed (182, 184, 186, 190, 192, 194, 196,
% 198, 200 in full; 181, 188, 189 are covers whose words become headings).
% Absent: pages 183, 185, 187, leaves of an unrelated English typescript
% (pp. 73, 72, 70 of a text on the cohomology H^i(V, O_V(nE)) of a variety,
% equations (15)-(18)), nothing in his hand — pages 182, 184, 186 are written
% on their versos; pages 191, 193, 195, 197, 199, leaves of a Bourbaki
% typescript (« Tribu n° 99 », « Algèbres de Lie semi-simples (5)-(9) »,
% scanned upside down), nothing in his hand — pages 190-200 (even) are their
% versos; page 181, an otherwise blank sheet inscribed top right in his hand
% « Généralités fonctorielles-catégoriques sur la distribution des Frobenius
% dans les gpes de Galois motiviques. », a cover, whose words become the
% first \section; page 188, a buff card in pencil, « Théorie de Galois des
% Motifs » (underlined) over « Généralités » (underlined), and page 189, a
% sheet inscribed top right « Foncteurs fibres motiviques et gp fondamental
% motivique. », two covers, whose words become the second \section and its
% \subsection*. None of them gets a \page.
%
% What the batch is. Two runs in his hand, neither paginated by him.
%   (1) Pages 182-186, after the cover of 181, a fast blue-ink draft,
%   « distribution des Frobenius »: M a category of geometric motives of
%   finite type over S of finite type over Z, degrees assumed even (else pass
%   to M^(pair)); closed points are « démultipliés » into pairs s' = (s, ξ),
%   i : G_m,C → G(s)_C; each « R-point motivique » s' gives a class
%   ξ(s') ∈ H^1(R, G^{0u}_R) (G^{0u} the kernel of ε : G^0_R → G_m) and a
%   homomorphism u(s') : G(s')_R → G_R^{ξ(s')} compatible with ε and i,
%   determined up to an element of G^{0u}_R(R). 184: the privileged
%   « π-structure » on G^0_R, i : μ_2,C → G^0_C, compared with
%   ψ(s') = u(s')_C ∘ (i(s')|μ_2); the two must agree, a non-tautological
%   condition on u(s'); with Tate's conjectures, knowing the u(s') amounts to
%   the effectivity structure on the G_R-modules. 186: the Frobenius
%   f(s') ∈ G_R^{ξ(s')}(R) have eigenvalues that are algebraic integers,
%   « une condition arithmétique assez draconienne ! ». The connective prose
%   of 182 and 184 is largely \ill{}.
%   (2) Pages 190-200, after the covers of 188-189, a steadier hand: the
%   canonical points of the motivic fundamental pro-torsor (T_α) of a
%   connected prescheme X, case by case. 1°) X of char. p > 0: no sections
%   in general; only ξ_ℓ ∈ T_α(Z_ℓ) (ℓ ≠ p, from the henselisation),
%   ξ_p ∈ T_α(W_p(k)) for k perfect (crystalline, with F, V, FV = p, ½-linear),
%   ξ_p^0 ∈ T_α(k) (« De Rham-Hodge »); N.B. obstructions in H^2 over F_q;
%   uniqueness up to isomorphism by Lang's theorem, with a guess that
%   T_∞(k) = lim T_α(k) has a single element (margin: « cf. Manin »). 192:
%   dependence on geometric points, π_1(X) acting on ξ_ℓ but not on ξ_p,
%   Aut(k/k(x)) acting ½-linearly; notations ξ_{x̄,ℓ}, ξ_{x,p}, ξ^0_{x,p}.
%   194-196: 2°) X dominating Spec Z: (i) ξ_ℓ for all ℓ, (ii) ξ_{x,p} at
%   points of char. p, with a conjectured relation between ξ_{y,p} and
%   ξ_{x,p} ⊗ W(k) (« ils ne sont en général pas isomorphes, sans doute »),
%   (iii) at a point with values in a field k of char. 0, ξ_{x,∞} ∈ T_α(k)
%   by Hodge-de Rham with its two filtrations (long vertical margin on the
%   canonical filtration of motives); for k = Q a trivialisation of T_Q,
%   no privileged isomorphism with ξ_{x̄,ℓ} ⊗ Q_ℓ, so every rational point
%   defines a class in H^1(Q_ℓ, G_ℓ^{(α)}). 198: (iv) at a complex point,
%   ξ_{x,Z} ∈ T_α(Z) (the integral lattice in F_x), ξ_{x,Z} ⊗ Z_ℓ ≅ ξ_{x,ℓ},
%   ξ_{x,Z} ⊗ C ≅ ξ_{x,∞}, the « involutive » isomorphism ξ_{x,Z} ≅ ξ_{τx,Z};
%   dependence on an archimedean place a) k(a) dense in C, b) k(a) real;
%   200: the real case, Gal(C/R) acting on ξ_{x,Z}, the invariants
%   ξ_{a,∞} ∈ T_α(R), ξ_{ā,Z} ∈ T_α(Z), ξ_{ā,∞} ∈ T_α(C) with
%   ξ_{ā,∞} ≅ ξ_{a,∞} ⊗_R C ≅ ξ_{ā,Z} ⊗_Z C; Aut_{k(a)}(C) acts ½-linearly.
%
% Continuations. Page 186 ends its run (the rest of the leaf is blank). Page
% 200 ends at the foot of the leaf with a complete sentence; whether the run
% continues on page 202 (batch 11) is not visible from here. The run of
% 190-200 continues across its own leaves: 194 ends « Il n'y a pas lieu de »
% and 196 opens « penser que … »; 196 ends « il faut remplacer Q_ℓ » and 198
% opens « par Q_ℓ(k) »; 198 ends « b) Si k(a) … dense dans » and 200 opens
% « C, i.e. … ». Nothing continues from batch 9, whose page 178 stops on a
% colon of a different run.
%
% Notation fixed for the batch. ξ with indices as he writes them (ξ_ℓ,
% ξ_ℓ^{(α)}, ξ_{x,p}, ξ_{x,p}^0, ξ_{x,∞}, ξ_{x,Z}, ξ_{ā,∞}); T_α for the
% components of the fundamental pro-torsor, \mathbf{T} for the bold T he
% writes a few times; Z̃_{ℓZ} (\widetilde{\mathbb{Z}_{\ell\mathbb{Z}}}) for
% the henselised local ring he writes with a tilde and glosses
% « hensélisé »; « ½ linéaire » for semi-linear, kept as he writes it
% (hand.md §2); G^{0u} for a superscript read « 0u », uncertain. « tt pt
% rat. », « pt géom. », « isom. can. », « val. », « gpe » kept.
% Plain square brackets in the text are his: words added in the interline.
%
% Legibility. Pages 190-200 are written fairly and read mostly through, the
% vertical margins of 192 and 194 less so. Pages 182 and 184 are a fast
% draft over a typescript show-through, and their connective prose is often
% \ill{}; the formulas carry.
%
% The dating is the inventory's own for the folder, copied verbatim. Its
% group, [10-18] « Théorie des motifs », is dated [à partir de 1958-à partir
% de 1982].
\documentclass[11pt,a4paper]{article}
% Le préambule commun à toutes les transcriptions du fonds Grothendieck.
%
% Il tient en une idée : **l'incertitude se compose, elle ne se résout pas.**
% Une transcription qui lisse un mot douteux a détruit la seule chose pour
% laquelle on la faisait — un lecteur doit pouvoir savoir, à chaque endroit,
% ce qui était sur le feuillet et ce qui a été ajouté. Les six macros qui
% suivent sont donc l'appareil critique, et non de la décoration.
%
% Les mêmes commandes sont reconnues par `scripts/render.mjs`, qui produit la
% vue de lecture du site. Une macro ajoutée ici doit l'être aussi là-bas,
% faute de quoi le rendu échouera bruyamment — ce qui est voulu : un rendu
% silencieusement faux est pire qu'un rendu absent.

\ProvidesPackage{grothendieck}[2026/08/08 Transcriptions du fonds Grothendieck]

\usepackage{amsmath,amssymb,amsthm}
\usepackage{mathrsfs}
\usepackage{stmaryrd}
% Les diagrammes commutatifs sont partout dans ce fonds : c'est la langue
% même de la géométrie algébrique de Grothendieck, et une transcription qui
% les rendrait en prose ne transcrirait rien.
\usepackage{tikz-cd}
% Français par défaut — c'est la langue des feuillets — mais l'anglais est
% chargé aussi : la traduction et le résumé sont des documents anglais, et la
% typographie française (espaces avant les deux-points) y serait une faute.
% Ces documents basculent par \selectlanguage{english} en tête de corps.
\usepackage[english,french]{babel}
\usepackage{fontspec}
\usepackage{geometry}
\geometry{a4paper,margin=25mm,marginparwidth=28mm}
\usepackage{marginnote}
\usepackage{xcolor}
% ulem, not soul. `soul`'s \st and \ul are built on a character-by-character
% scan that XeLaTeX's font machinery breaks: the document fails with an
% undefined control sequence rather than degrading. ulem does the same two jobs
% and is Unicode-engine safe. `normalem` stops it hijacking \emph, which the
% transcriptions use for Grothendieck's own emphasis.
\usepackage[normalem]{ulem}
\usepackage{hyperref}
\hypersetup{colorlinks=true,linkcolor=black,urlcolor=black,pdfborder={0 0 0}}

% --- Métadonnées du lot -----------------------------------------------------
% Reprises telles quelles par le script de rendu, qui les affiche en tête de
% la vue de lecture. Elles fixent ce que le fichier prétend couvrir : sans
% elles, une transcription flotte sans point d'attache dans un fonds de seize
% mille feuillets.

\newcommand{\folder}[1]{\gdef\grothFolder{#1}}
\newcommand{\batch}[1]{\gdef\grothBatch{#1}}
\newcommand{\leaves}[2]{\gdef\grothFirst{#1}\gdef\grothLast{#2}}
\newcommand{\foldertitle}[1]{\gdef\grothFoldertitle{#1}}
\gdef\grothFolder{?}\gdef\grothBatch{?}\gdef\grothFirst{?}\gdef\grothLast{?}\gdef\grothFoldertitle{}

% --- Le résumé --------------------------------------------------------------
%
% Il ouvre la lecture modernisée, sous le titre « Résumé », et il est composé
% en petit. Ce n'est pas un ornement : c'est le seul passage du document qui
% ne soit pas des mathématiques, et un lecteur doit pouvoir le reconnaître
% avant de l'avoir lu — puis le sauter s'il connaît déjà le sujet, ou s'y
% arrêter s'il ne le connaît pas. Le filet qui le suit marque la reprise du
% corps.

\newenvironment{resume}
  {\small\setlength{\parskip}{0.45em}}
  {\par\medskip\hrule\medskip}

% --- Mots-clés ---------------------------------------------------------------
%
% En anglais, en clôture du résumé de la lecture modernisée : le vocabulaire
% moderne sous lequel chercher ce que les feuillets construisent. Le manifeste
% du site (`npm run manifest`) les extrait pour étiqueter la cote entière —
% cette ligne est la source unique des tags, il n'y a pas de fichier à côté.
\newcommand{\keywords}[1]{\par\smallskip\noindent
  {\footnotesize\textsc{Keywords} --- \textit{#1}}\par}

% --- L'appareil critique ----------------------------------------------------

% Le numéro de feuillet, en marge. C'est l'ancre qui permet de régler un
% désaccord sur une formule en regardant *un* feuillet plutôt qu'une chemise
% de six cents. Le site s'en sert aussi pour tourner les pages du fac-similé
% quand on fait défiler la transcription.
\newcommand{\leaf}[1]{%
  \marginnote{\footnotesize\color{gray!70}\textsf{#1}}%
  \label{leaf:#1}\ignorespaces}

% Illisible : on ne devine pas. Un mot inventé qui se lit comme les autres est
% le pire résultat possible.
\newcommand{\ill}{\textcolor{red!60!black}{[\,\ldots\,]}}

% Lecture douteuse : le mot est donné, et signalé comme incertain.
\newcommand{\uncertain}[1]{\uline{#1}}

% Ajout de l'éditeur — une lettre manquante, un mot sous-entendu.
\newcommand{\add}[1]{\textcolor{blue!50!black}{[#1]}}

% Biffé par Grothendieck lui-même. Ce qu'il a rayé fait partie du document :
% une rature dit souvent où la pensée a changé de direction.
\newcommand{\struck}[1]{\sout{#1}}

% Note du transcripteur, sur l'état matériel du feuillet.
\newcommand{\note}[1]{\par\smallskip{\small\color{gray!60!black}\textit{[#1]}}\par\smallskip}

% Note marginale de Grothendieck, distincte du corps du texte.
\newcommand{\marginal}[1]{\par\smallskip\noindent\hspace{1em}%
  {\small\color{gray!60!black}#1}\par\smallskip}

% --- Titre ------------------------------------------------------------------

\AtBeginDocument{%
  \begin{center}
    {\large\bfseries Fonds Alexandre Grothendieck --- cote n\textsuperscript{o}~\grothFolder}\\[2pt]
    {\itshape\grothFoldertitle}\\[4pt]
    {\small feuillets \grothFirst--\grothLast\ (lot \grothBatch)}\\[2pt]
    {\footnotesize\color{gray!60!black}Transcription établie d'après le fac-similé numérisé
     par l'Université de Montpellier.\\
     Les lectures incertaines sont soulignées, les illisibles notées [\,\ldots\,],
     les ajouts entre crochets.}
  \end{center}
  \vspace{1em}}

\endinput


\folder{15}
\batch{10}
\pages{181}{200}
\dating{1965-[vers 1977]}
\watermark{Édition de démonstration}
\foldertitle{Théorie arithmétique. Théorie de Galois des motifs (1965) : notes manuscites (s.d.), tapuscrit (s.d.).}

\begin{document}

\section{Généralités fonctorielles-catégoriques sur la distribution des
Frobenius dans les gpes de Galois motiviques}

\note{ces mots sont de sa main, seuls, en haut à droite du feuillet 181, qui
sert de couverture ; ils deviennent ici le titre de section}

\page{182}
Soit $M$ une catégorie de motifs géométriques, [de type fini,] sur le
préschéma $S$ de type fini sur $\mathbb{Z}$. Pour simplifier, supposons que
les degrés de $M$ soient pairs (sinon, il sera possible, pour
[\uncertain{néanmoins}] \struck{\ill{}} [traiter] la question de la
distribution des Frobenius, de remplacer $M$ par $M^{(\mathrm{pair})}$).
Tout \uncertain{choix} d'un \ill{} permet de définir des pts fermés
$s \in S$ \ill{} « $\mathbb{R}$-points motiviques » \struck{de $M$}
\ill{} \ill{}. Pour fixer les idées, \ill{} \ill{} un pt
$\mathbb{R}$-\uncertain{valué} \ill{} \ill{} \uncertain{algébrique}
$G_{\mathbb{R}}$ sur $\mathbb{R}$ de $M$, \ill{} \ill{} \ill{}.

Il y a lieu de \struck{\ill{}} [\uncertain{démultiplier}] les points
[fermés] de $S$, en considérant les couples $s' = (s, \xi)$, où
$s \in S^{0}$, et \ill{} \struck{\ill{}} \uncertain{géométrique} \ill{}
est du type
\[
  i : \mathbb{G}_{m,\mathbb{C}} \to G(s)_{\mathbb{C}}
\]
\ill{} Galois motivique associé à $M_s$, --- $\mathbb{C}(s)$ désigne la
\ill{} \ill{} \ill{} \uncertain{complexe}. À tout $s'$ \ill{}

\marginal{l'existence d'un \ill{} \uncertain{désigne} \ill{}
\uncertain{généraux}}
\note{note oblique dans la marge gauche, le dernier mot souligné}

\struck{Alors} Soit $S'$ l'ensemble des $\mathbb{R}$-points motiviques de
$M$, \struck{et} \ill{} \ill{} déterminé à isom.\ près \struck{\ill{}}
\[
  \xi(s') \in H^{1}(\mathbb{R}, G^{0u}_{\mathbb{R}})
\]
où $G^{0u}$ est le noyau de $G^{0}_{\mathbb{R}}
\xrightarrow{\ \varepsilon\ } \mathbb{G}_m$. De plus, si
$G_{\mathbb{R}}^{\xi(s')}$ est le « \uncertain{tordu} » \ill{} de
$G_{\mathbb{R}}$ par torsion par $\xi(s')$, on obtient \ill{} un
homomorphisme \ill{}
\[
  u(s') : G(s')_{\mathbb{R}} = G(s)_{\mathbb{R}} \longrightarrow
  G_{\mathbb{R}}^{\xi(s')}
\]
compatible avec $\varepsilon_{\mathbb{R}}$, $i_{\mathbb{R}}$, et déterminé
en fait mod \struck{automorphismes intérieurs} [\ill{}] par un élément de
$G^{0u}_{\mathbb{R}}(\mathbb{R})$.
\note{l'exposant de $G^{0u}$ se lit mal : un rond suivi d'une boucle, lu
« $0u$ » ; on pourrait aussi y voir « $0n$ ». Un premier $G$ en tête de la
formule $\xi(s')$ est surchargé}

\page{184}
De plus, on a une structure [$\pi$] privilégiée sur $G^{0}_{\mathbb{R}}$,
qui peut \struck{définir} s'exprimer définie par un hom
\struck{\ill{}}
\[
  i : \mu_{2,\mathbb{C}} \longrightarrow G^{0}_{\mathbb{C}}
\]
(comme on \ill{} en envoyant un tore [\ill{}] (\ill{})). D'autre part,
$G(s')_{\mathbb{C}}$ est muni de $i(s') : \mathbb{G}_{m,\mathbb{C}} \to
G(s')_{\mathbb{C}}$, et par restriction [à $\mu_2$] on obtient
canoniquement
\[
  \psi(s') = u(s')_{\mathbb{C}} \circ \bigl(i(s') \vert \mu_2\bigr) :
  \mu_{2,\mathbb{C}} \longrightarrow G_{\mathbb{C}}^{\xi(s')}
\]
\struck{D'autre part} Ceci posé, il faut que $\pi^{\xi(s')}$ (\ill{} de
$G_{\mathbb{R}}^{\xi(s')}$ déduite de $\pi$ sur $G$ par transport) soit
celle définie par les $\mu_{2,\mathbb{C}}$.
\note{la lettre du premier membre, $\psi(s')$, est douteuse : une sorte de
$\mathrm{v}$ à queue}

Par conséquent, puisqu'il s'agit [de \ill{}] de $G_{\mathbb{R}}^{\xi(s')}$
[\ill{}] que $\xi(s')$ \ill{} [\uncertain{à conjugaison près}], les
$\pi$-structures en question doivent être identiques --- \ill{} $u(s')$
doivent être une condition non tautologique [\ill{}] ce qui donne une
[\uncertain{restriction}] sur $u(s')$). [Notons cependant que la
distribution des $\pi$-structures d'un gpe réel $H$, définie en termes des
hom.\ $\mu_{2,\mathbb{C}} \to H_{\mathbb{C}}$, \ill{} \ill{}
\uncertain{commutatif} en tous cas !]
\note{« identiques » et « tautologique » portent un trait au-dessus ; dans
la parenthèse finale, un mot est biffé avant « $\pi$-structures »}

? d'\uncertain{ailleurs} \ill{} \uncertain{finie}, [si \ill{}] \ill{}
\ill{} \ill{} \uncertain{dominance}. \ill{} les \uncertain{restrictions}
\ill{} \ill{}
\note{deux lignes rapides, précédées d'un « ? » dans la marge gauche, dont
presque rien ne se lit}

Enfin, conjecturalement \ill{} [\uncertain{avec les conjectures} de]
Tate vérifiées, la connaissance des $u(s')$ \struck{\ill{}} \ill{} de la
structure d'effectivité dans la catégorie

\page{186}
des $G_{\mathbb{R}}$-modules de dim.\ finie sur $\mathbb{R}$. La
\uncertain{condition} qui exprime cette effectivité est \struck{singulière}
par les \struck{réels} noyaux définis par les $v(s')$ [sous forme
« \uncertain{analytique} »]. C'est un fait qui exprime [le fait que les]
Frobenius
\[
  f(s') \in G_{\mathbb{R}}^{\xi(s')}(\mathbb{R})
\]
ont \ill{} par leurs valeurs propres, dans \ill{} représentation \ill{} de
$G_{\mathbb{R}}^{\xi(s')}$ \struck{$(\mathbb{R})$} \ill{} \ill{} de
$G_{\mathbb{C}}^{\xi(s')}$ --- sont des \emph{entiers} algébriques, ce qui
est une condition arithmétique assez draconienne !
\note{« entiers » est souligné. Le premier mot de la deuxième ligne,
« singulière », est souligné puis biffé ; l'exposant de $f$ se lit
$(s')$. Le texte s'arrête au tiers supérieur du feuillet, le reste est
blanc}

\section{Théorie de Galois des Motifs. Généralités}

\note{deux lignes de sa main, au crayon, soulignées, seules sur le feuillet
188, qui sert de couverture ; elles deviennent ici le titre de section}

\subsection*{Foncteurs fibres motiviques et gp fondamental motivique}

\note{titre de sa main, seul sur le feuillet 189, qui sert de couverture}

\page{190}
$X$ préschéma connexe, \struck{réduit, ce qui ne restreint pas}

1°) $X_{\mathrm{réd}}$ de car.\ $p > 0$. Alors dans le pro-torseur
fondamental $(T_\alpha)$, il n'y a \uncertain{pas} en général de sections
pour les $T_\alpha$.

On dispose uniquement, \struck{de \uncertain{façon} canonique}
\marginal{Coh.\ $\ell$-adique}
\begin{enumerate}
  \item[(i)] des $\xi_\ell^{(\alpha)} \in T_\alpha(\mathbb{Z}_\ell)$,
  $\ell \neq p$, provenant sans doute de $\xi_\ell \in
  T_\alpha(\widetilde{\mathbb{Z}_{\ell\mathbb{Z}}})$
  \marginal{hensélisé}
  \item[(ii)] d'un $\xi_p^{(\alpha)} \in T_\alpha(W_p(k))$ si
  \struck{$X$ déf.\ \ill{}} $k$ est un corps parfait sur $X$
  \struck{le corps parfait $k$}
  [$\xi_p$ provenant sans doute d'un $\xi_p \in
  T_p(\widetilde{\mathbb{Z}_{p\mathbb{Z}}})$]
  \marginal{hensélisé strict.}
  \item[(iii)] \uncertain{réduit} $\xi_p^{(\alpha)\,0} \in T_\alpha(k)$
  \marginal{« De Rham-Hodge »}
\end{enumerate}
\note{les items (ii) et (iii) sont réunis par une accolade, précédée d'un
« ? ». « hensélisé » et « hensélisé strict » sont écrits sous les deux
anneaux marqués d'un tilde, qu'ils glosent. « De Rham-Hodge » est écrit
au-dessus de (iii), avec une flèche vers le $k$ de $T_\alpha(k)$}

\marginal{Coh.\ $p$-adique. Structure supplémentaire : opérations $F$ et
$V$, ½ linéaires \uncertain{relativement} au Frobenius \ill{}}
\marginal{Coh.\ Hodge-de Rham. Mais ici, il faut \struck{sans doute}
travailler dans la catégorie \uncertain{dérivée}. Tenir compte aussi de la
\uncertain{filtration} Hodge-De Rham}
\note{ces notes marginales sont écrites en oblique dans la marge gauche, en
face de (i), (ii), (iii) respectivement : la deuxième est entourée d'un trait
qui la rattache à (ii) ; dans la deuxième, « ½ linéaires relativement au
Frobenius » est surchargé de plusieurs traits}

\uncertain{N.B.} Si $X = \operatorname{Spec} \mathbb{F}_q$, il y a en
général des \struck{\ill{}} obstructions non triviales dans des
$H^2(\mathbb{Q}, G)$ ---
\note{le $\mathbb{Q}$ de $H^2(\mathbb{Q}, G)$ est douteux ; le signe
initial, en marge, est une sorte de monogramme, lu N.B.}

On utilise des « foncteurs fibres » qui induisent sur les motifs effectifs
de poids $0$ le foncteur canonique. Alors les \uncertain{points} de
$T_\alpha(\mathbb{Z}_\ell)$ sont sans doute uniques [à isom.\ près] (th.\ de
Lang), de même ceux de $T_\alpha(W_p(k))$ si $k$ fini (ii), [plus précis.]
pour un $k$ parfait de dim.\ coh.\ $1$ (?). En général, il faut s'attendre
que $T_\alpha(k)$ \struck{\ill{}} ait plus d'un élément à isom.\ près, mais
il est possible que $T_\infty(k) = \varprojlim T_\alpha(k)$ n'ait qu'un
élément à isom.\ près, celui donné par la théorie de De Rham-Hodge.

\marginal{\uncertain{relation} \ill{} \uncertain{galoisienne} ?? Ce doit
être équivalent, cf.\ \emph{Manin}}
\note{note écrite en oblique dans la marge gauche, en bas, reliée par un
long trait de crayon à la note sur la cohomologie $p$-adique ; « Manin » est
souligné}

\page{192}
les $\xi_\ell$ \ill{} dépendant du choix d'un pt géom.\ de $X$, mais à
isom.\ unique près (dépendance du $\pi_1(X)$) [$\pi_1(X)$ opère sur
$\xi_\ell$] \struck{sont} indépendants de ce choix.
\note{« dépendance du $\pi_1(X)$ » est souligné deux fois}

Les $\xi_p^{\alpha}$ dépendent du choix d'un \struck{corps} pt géom.\ [de
$X$] à valeurs dans un corps parfait sans plus, si on veut les
$k(x)^{\mathrm{pf}}$ (clôture parfaite de $k(x)$), et pour $k$ donné, il ne
semble pas qu'il y ait de choix des isom.\ entre les $\xi_{p,a}^{(\alpha)}$
correspondant aux divers pts à valeurs dans $k$ (c'est-à-dire « points
\ill{} » du \ill{} géom.\ qui \ill{} des \uncertain{maximaux}), s'ils sont
localisés en le même $x \in X$, \struck{\ill{}} \struck{\ill{}} si $k$ est
alg.\ clos, ils sont conjugués par autom.\ de $k$ [(éléments de
$W_p(k)$)] \ill{} alg., mais sont \ill{} conjugués \ill{} autom.\ de $k$
… il est à \ill{} général qu'il semble \ill{} l'un \ill{} de
l'autre… Le gpe $\pi_1(X)$ n'opère plus sur $\xi_p^{\alpha}$,
\marginal{En revanche, on a comme structure supplémentaire les opérations
$F$, $V$ avec $FV = p$}
par contre le groupe des autom.\ de $k$ sur $k(x)$
\struck{\ill{}} opère ½ linéairement sur $\xi_p^{(\alpha)}$ ---
\ill{} fibre $\ell$-adique en \ill{}
\note{« ½ linéairement » est ajouté dans l'interligne et souligné ; une
insertion biffée au-dessus de la ligne se termine par un signe de renvoi qui
mène à la note marginale ci-dessous}

\marginal{\ill{} élégante. D'autre part, un autom.\ $\sigma$ de $k$ qui
induit sur $k(x)$ l'identité transforme en général $\xi_{x,p}$ en un
$\xi_{x,p}^{\sigma}$ qui n'est plus isom.\ à $\xi_{x,p}$ ---}
\note{les deux notes marginales sont écrites verticalement dans la marge
gauche ; dans la seconde, les indices et exposants des deux premiers $\xi$
sont douteux}

Notations,
\begin{enumerate}
  \item[(i)] $\xi_{\bar{x},\ell}$ \ill{} ; $\xi_{x,\ell}$ \ill{} pour
  $x \in X$ ou un pt de $x$ à val.\ dans un corps quelc.\ ($\xi_{x,\ell}$
  déf.\ à isom.\ non unique près)
  \item[(ii)] $\xi_{x,p}$, $\xi_{x,p}^{0}$ : $x$ pt de $X$ à val.\ dans un
  corps \uncertain{parfait}.
\end{enumerate}

\marginal{\uncertain{Gros} point \uncertain{subtil} \ill{} \ill{}}
\note{note de trois lignes soulignées, en bas de la marge gauche}

\page{194}
2°) $X$ \struck{réd} dominant $\operatorname{Spec} \mathbb{Z}$, i.e.\ $X$
contient des pts de car.\ résiduelle nulle.

Alors \struck{en plus \ill{} [dessus] \ill{} éventuellement \ill{} (iii)
\ill{} d'inégales car.) \ill{}}
\struck{(iii bis)}
\note{deux lignes biffées d'un trait}

\begin{enumerate}
  \item[(i)] $\xi_\ell \in \mathbf{T}(\mathbb{Z}_\ell)$ $\forall\, \ell \in
  \mathbb{P}$, $\xi_\ell$ provenant sans doute de
  $T_\alpha(\widetilde{\mathbb{Z}_{\ell\mathbb{Z}}})$.

  $\xi_\ell$ dépend du choix d'un pt géom.\ quelconque de $X$, et est
  indép.\ de ce choix modulo opération de $\pi_1$, notation
  $\xi_{\bar{x},\ell}$, $\xi_{x,\ell}$.

  \item[(ii)] Si $X$ a des points $x$ de car.\ $p > 0$, on a \ill{}
  également des $\xi_{x,p}$, $\xi_{x,p}^{0}$.

  Conjecture : relation, \struck{\ill{}} si $y$ est un pt de car.\ $p > 0$
  et \struck{\ill{}} un pt de car.\ $0$, entre $\xi_{y,p}$ et
  $\xi_{x,p} \otimes_{\mathbb{Z}_p} W(k)$ [$y$ pt à val.\ dans un corps
  parfait $k$]

  Ils ne sont en général pas isomorphes, sans doute.
  \struck{\ill{} ils sont \ill{} \ill{}}
  \note{« Ils ne sont en général pas isomorphes » est souligné ; suivent
  deux lignes surchargées et biffées, que la marge gauche désigne d'un mot
  lui-même biffé}

  \item[(iii)] Si $x$ est un pt à val.\ dans un corps [$k$] de car.\ $0$,
  on a $\xi_{x,\infty} \in T_\alpha(k)$ par Hodge-de Rham [avec ses
  \uncertain{2} filtrations \struck{\ill{}}]. Il n'y a pas lieu de
\end{enumerate}
\note{« car.\ 0 » est souligné ; l'insertion « avec ses 2 filtrations » se
rattache à « Hodge-de Rham », et le « 2 » est entouré}

\marginal{l'une est celle qui donne comme gradué associé un
$H^{\bullet}(X, \Omega^{P})$, l'autre (\uncertain{authentique}) fait penser
à : la filtration canonique des motifs. (Celle-ci est définie pour les
motifs \ill{}, \ill{} filtration est \uncertain{originaire} du foncteur
Hodge-De Rham, bien qu'on s'attende qu'elle ne provienne pas d'une
filtration sur \ill{} motif, c'est une filtration en \ill{}…}
\note{note écrite verticalement dans la marge gauche, en face de (iii),
rattachée par un trait aux « 2 filtrations »}

\page{196}
penser que pour ces pts définis à val.\ dans le même $k$, les deux
foncteurs fibres, soient isom.\ (pas plus que dans le cas (ii)). Il n'y a
plus en général d'opération de $\pi_1(X)$ sur $\xi_{x,\infty}$, mais encore
comme dans (ii) les autom.\ de $k$ sur $k(x)$ opèrent sur $\xi_{x,\infty}$
½ linéairement (donnée de descente). Un autom.\ de $k$ qui \ill{}
l'identité sur $k$ définit $\xi_{x,\infty}^{\sigma}$ qui à priori n'est plus
isomorphe à $\xi_{x,\infty}$.
\note{« ½ linéairement » est souligné}

Un cas intéressant est celui où $k = \mathbb{Q}$, on trouve alors un
\struck{trivial} $\xi_{x,\infty} \in T_\alpha(\mathbb{Q})$, i.e.\ une
trivialisation de $T_{\mathbb{Q}}$. Deux trivialisations \ill{} de
$T_{\mathbb{Q}}$ ne sont en général pas \uncertain{isom.} à priori
\uncertain{les mêmes}. De plus, il n'y a pas d'isom.\ canoniques
privilégiés entre $\xi_{x,\infty}$ \ill{} $\xi_{x,\infty}
\otimes_{\mathbb{Q}} \mathbb{Q}_\ell$ et $\xi_{\bar{x},\ell}
\otimes_{\mathbb{Z}_\ell} \mathbb{Q}_\ell$ --- \struck{bien que pas} et
comme la première à priori ne provient pas d'un objet de
$T_\alpha(\mathbb{Z}_\ell)$ mais seulement de $T(\mathbb{Q}_\ell)$, il n'y a
pas lieu de penser qu'il y ait un isom.
\marginal{en effet, ils \uncertain{sont} \uncertain{divisés} \ill{} sur
les \uncertain{motifs} de \uncertain{poids et niveau} $0$}
Ainsi, tt pt rat.\ $/\mathbb{Q}$ définit une classe de coh.\ dans les
$H^{1}(\mathbb{Q}_\ell, G_\ell^{(\alpha)})$. Dès \ill{} \struck{\ill{}} en
général, il faut remplacer $\mathbb{Q}_\ell$
\note{« canoniques » est souligné. La note marginale, oblique, est
séparée du texte par un trait vertical en face de « ne provient pas d'un
objet de $T_\alpha(\mathbb{Z}_\ell)$ ». L'exposant de $H^{1}$ peut se lire
1 ou 2 (comparer le $H^2$ du feuillet 190) ; le $G$ porte un rond
au-dessus, peut-être un $G^{0}$}

\page{198}
par $\mathbb{Q}_\ell(k)$ --- c'est particulièrement intéressant dans le cas
\ill{} un corps de nbs.

\begin{enumerate}
  \item[(iv)] Si $x$ est un pt à valeurs dans $\mathbb{C}$, il lui
  correspond
  \[
    \xi_{x,\mathbb{Z}} \in \mathbf{T}_\alpha(\mathbb{Z}) \qquad
    \xi_\alpha(\mathbb{Z})(F) = \text{réseau entier dans } F_x ,
  \]
  \struck{Quand} \struck{De plus, le gpe \ill{} $\xi_{x,\mathbb{Z}}$ et
  $\xi_{x,\mathbb{Z}}$ \ill{} $\mathbb{Z}/2\mathbb{Z}$ opère \ill{}}
  \struck{\ill{} complexe)}, et des isom.\ can.
  \[
    \begin{cases}
      \xi_{x,\mathbb{Z}} \otimes_{\mathbb{Z}} \mathbb{Z}_\ell \simeq
      \xi_{x,\ell} \\
      \xi_{x,\mathbb{Z}} \otimes_{\mathbb{Z}} \mathbb{C} \simeq
      \xi_{x,\infty}
    \end{cases}
  \]
\end{enumerate}
\note{« dans $\mathbb{C}$ » est souligné. Entre les deux formules, deux
lignes et un encadré sont biffés ; on y lit $\xi_{x,\mathbb{Z}}$,
$\mathbb{Z}/2\mathbb{Z}$ et « complexe »}

D'ailleurs $\xi_{x,\mathbb{Z}}$ est défini en termes d'un pt à val.\ dans
un corps [top.] $K$ isom.\ à $\mathbb{C}$ (i.e.\ isom.\ \ill{}) en d'autres
termes il y a un isom.\ canonique « involutif »
\[
  \xi_{x,\mathbb{Z}} \simeq \xi_{\tau x,\mathbb{Z}}
\]
où $\tau x$ est le pt complexe conjugué.

\struck{Donc} On peut dire aussi si $x$ est localisé en $a \in X$, que [a)
si $k(a)$ est dense dans $\mathbb{C}$, alors,] $\xi_{x,\mathbb{Z}}$ dépend
de la donnée d'une topologie [archimédienne] (ou valuation) sur $k(a)$, à
complété $\mathbb{C}$ \quad b) Si $k(a)$ \ill{} dense dans

\page{200}
$\mathbb{C}$, i.e.\ $\overline{k(a)} \simeq \mathbb{R}$ [isom.\ unique
can.], alors $\xi_{x,\mathbb{Z}}$ est encore défini par la valuation
archimédienne sur $k(a)$, mais \struck{suivant} [compte tenu] \struck{des}
les choix possibles de la clôture algébrique de $k(a)$, le gpe de Galois de
$\mathbb{C} = \overline{\mathbb{R}}$ sur $\mathbb{R}$ opère dans
$\xi_{x,\mathbb{Z}}$, \ill{} [les invts \ill{} $\xi_{x,\mathbb{Z}}$ \ill{}
de $\xi_{x,\mathbb{Z}}$, \ill{}] [\ill{} \ill{}] et cette opération, qui
est compatible avec \ill{} l'action ½ linéaire sur $\xi_{\bar{a},\infty}$
correspondant à $\tau$, il n'y est pas question par contre d'un
\uncertain{réseau stable}].
\note{« et cette opération » est souligné deux fois, « ½ linéaire »
une fois}

\marginal{\ill{} a choisi \ill{} \uncertain{clôture algébrique} \ill{} à
val.\ dans $\mathbb{C}$}

Dans ce cas, les invariants sont

\[
  \begin{array}{l}
    \xi_{a,\infty} \in T_\alpha(\mathbb{R}) \\
    \xi_{\bar{a},\mathbb{Z}} \in T_\alpha(\mathbb{Z})
  \end{array}
  \qquad
  \xi_{\bar{a},\infty} \in T_\alpha(\mathbb{C})
\]
\note{sur la page, les deux éléments de gauche sont reliés à celui de droite
par deux traits obliques ; au-dessus de $\xi_{\bar{a},\mathbb{Z}}$, une
première formule est biffée, $\xi_{a,\infty} \otimes_{\mathbb{R}}
\mathbb{C} \simeq$}

\[
  \xi_{\bar{a},\infty} \overset{\mathrm{can}}{\simeq} \xi_{a,\infty}
  \otimes_{\mathbb{R}} \mathbb{C} \simeq \xi_{\bar{a},\mathbb{Z}}
  \otimes_{\mathbb{Z}} \mathbb{C}
\]

$\pi_1(X)$ n'opère plus à priori dans $\xi_{a,\mathbb{Z}}$, les seules
opérations sont celles de $\tau$ si $a$ est réel. Cependant, sur
$\xi_{\bar{a},\infty}$, il y a (en plus éventuellement de $\tau$, si $a$ est
réel) les opérations de \uncertain{$\mathrm{Aut}_{k(a)}(\mathbb{C})$}, qui
sont ½ linéaires.
\note{« réel » est souligné dans la première phrase}

\end{document}
